% ============================================================
%  TERM PROJECT – EXAMPLE DOCUMENTATION
%  Databases · THGA Bochum
%
%  This is a worked example of the *documentation* deliverable.
%  It shows the expected structure and level of detail. Your own
%  documentation describes your own application.
% ============================================================
\documentclass[11pt,a4paper]{article}

\usepackage{thga-db}

\renewcommand{\thgaKurs}{Databases}
\renewcommand{\thgaDatum}{Summer Term 2026}

\begin{document}
\thispagestyle{empty}
\thgaTitel{Pet Adoption Management System}%
{Club Library: a database-backed lending system}

\begin{infobox}
\textbf{Author:} Kaddoum Meriem \\
\textbf{Module:} Introduction to Database Management Systems · THGA Bochum \\
\textbf{Stack:} PostgreSQL · FastAPI · Docker Compose · tkinter (\texttt{.deb}) \\
\textbf{Repository:} \url{https://github.com/meryamkaddoum-max/pet-adoption-system}
\textbf{Video:} \url{https://www.youtube.com/watch?v=xrGPAW4uFLM}
\end{infobox}

\tableofcontents
\vspace{1em}

% ============================================================
\section{Introduction and motivation}
% ============================================================


Animal shelters and adoption centers often rely on manual processes or simple tools
such as spreadsheets to manage their data. While this may be sufficient for small
datasets, it becomes inefficient and error-prone when the amount of data increases
or when multiple entities and relationships must be handled simultaneously.

The \textbf{Pet Adoption Management System} was developed to address these limitations
by providing a structured and database-driven solution. The system manages three
core entities: pets, adopters, and adoptions. These entities are stored in a relational
database and can be accessed through a REST API.

The main goal of the system is to provide a consistent and reliable way to store,
retrieve, and manipulate data related to the adoption process. By using a backend
API, the system ensures that all data operations follow defined rules and constraints.

A graphical user interface (GUI) was implemented using tkinter to allow users to
interact with the system in a simple and intuitive way. The frontend communicates
with the backend exclusively via HTTP requests, ensuring a clear separation of
responsibilities between presentation, logic, and data storage.

\begin{beispielbox}
\textbf{Usage scenario.}
A user wants to adopt a pet. First, the system displays all available pets using the
API. The user selects a pet and provides adopter information. The adoption is then
recorded, linking the pet with the adopter. The system ensures that the data is
stored consistently in the database.
\end{beispielbox}
% ============================================================
\section{Data model}
% ============================================================


During the implementation of the system, the data model was extended in order to
better reflect real-world requirements. Additional attributes were introduced to
improve usability and data completeness.

\textbf{Extended attributes:}
\begin{itemize}
\item Pet: breed, gender, status (e.g. available, adopted)
\item Adopter: first\_name, last\_name, phone
\item Adoption: adoption\_date
\end{itemize}

The system is based on three main entities: \textbf{pets}, \textbf{adopters}, and
\textbf{adoptions}.

\textbf{Relationships:}
\begin{itemize}
\item One adopter can adopt multiple pets over time
\item Each pet can only be adopted once
\end{itemize}

This results in two one-to-many relationships:
\begin{itemize}
\item adopter (1) -- (N) adoption
\item pet (1) -- (N) adoption
\end{itemize}

The \textbf{adoption} entity acts as a linking table between pets and adopters.
It contains foreign keys referencing both entities and stores additional information
such as the adoption date.

\subsection{Relational schema}

Based on the conceptual model, the following relational schema was implemented
(primary keys are underlined, foreign keys are italicized):

\begin{itemize}
\item \texttt{adopter}(\underline{id}, first\_name, last\_name, email, phone)
\item \texttt{pet}(\underline{id}, name, species, breed, gender, age, status)
\item \texttt{adoption}(\underline{id}, \emph{adopter\_id}, \emph{pet\_id}, adoption\_date)
\end{itemize}

\textbf{Normalisation.}

All tables are designed according to the Third Normal Form (3NF). Each attribute
depends only on the primary key of its table, and no redundant data is stored.
Relationships between entities are represented using foreign keys, ensuring
referential integrity.

% ============================================================
\section{Database (PostgreSQL)}
% ============================================================


The database is implemented using PostgreSQL and is initialized automatically
when the Docker container starts. SQL scripts are used to define the schema and
insert initial test data.

\begin{lstlisting}[language=SQL, caption={schema.sql (excerpt)}]

CREATE TABLE pets (
    id SERIAL PRIMARY KEY,
    name VARCHAR(120),
    species VARCHAR(80),
    breed VARCHAR(120),
    gender VARCHAR(20),
    age INTEGER,
    status VARCHAR(20)
);

CREATE TABLE adopters (
    id SERIAL PRIMARY KEY,
    first_name VARCHAR(120),
    last_name VARCHAR(120),
    email VARCHAR(200),
    phone VARCHAR(50)
);

CREATE TABLE adoptions (
    id SERIAL PRIMARY KEY,
    adoption_date DATE,
    pet_id INTEGER REFERENCES pets(id),
    adopter_id INTEGER REFERENCES adopters(id)
);

\end{lstlisting}

The use of PostgreSQL ensures reliable data storage, transaction safety, and
support for relational constraints such as foreign keys.

% ============================================================
\section{REST API (FastAPI)}
% ============================================================


The backend of the system is implemented using FastAPI and acts as the central
interface between the database and the frontend.

The API provides endpoints to perform CRUD operations on the main entities.
All database interactions are handled through SQLAlchemy ORM, which maps Python
objects to database tables.

\textbf{Supported operations:}
\begin{itemize}
\item Create, read, update, and delete pets
\item Create, read, and delete adopters
\item Create, read, and delete adoptions
\end{itemize}

\textbf{Example endpoints:}

\begin{tabular}{lll}
GET & /pets & Retrieve all pets \\
POST & /pets & Create a new pet \\
DELETE & /pets/\{id\} & Delete a pet \\
POST & /adoptions & Create a new adoption \\
\end{tabular}

\textbf{Validation.}

Pydantic schemas are used to validate incoming data and ensure correct data types.
This prevents invalid data from being stored in the database.

The API follows a clear separation of concerns: routing defines endpoints, while
business logic and database access are handled separately.

% ============================================================
\section{Frontend (tkinter)}
% ============================================================


The frontend is implemented using Python tkinter and provides a graphical user
interface for interacting with the system.

The application consists of three main tabs:

\begin{itemize}
\item Pets: create, view, and delete pets
\item Adopters: manage adopter information
\item Adoptions: create and manage adoptions
\end{itemize}

The frontend communicates with the backend using HTTP requests via the Python
requests library. No direct database access is performed from the frontend.

\textbf{Architecture.}

The system follows a layered architecture:
\begin{itemize}
\item Presentation layer: tkinter GUI
\item Application layer: FastAPI backend
\item Data layer: PostgreSQL database
\end{itemize}

\textbf{Implementation challenges.}

A key challenge was ensuring correct handling of the \textbf{status attribute}
for pets. The frontend had to correctly send user input to the API and display
the returned values consistently.

Additionally, proper error handling was implemented to display meaningful
messages to the user in case of API or input errors.
% ============================================================
\section{Operation and deployment}
% ============================================================


The application is deployed using Docker Compose, which allows multiple services
to run together in isolated containers.

The system consists of two main services:
\begin{itemize}
\item PostgreSQL database
\item FastAPI backend
\end{itemize}

Docker ensures that the application runs in a consistent environment without
manual installation of dependencies.

\textbf{Deployment steps:}
\begin{enumerate}
\item Run \texttt{docker compose up --build}
\item Wait until the database and API are started
\item Access the API via \texttt{http://localhost:8000}
\end{enumerate}



% ============================================================
\section{Reflection}
% ============================================================


The project demonstrates the benefits of separating the system into distinct
layers: database, backend API, and frontend.

This modular architecture allowed each component to be developed and tested
independently. For example, the API could be tested using HTTP requests before
integrating it with the frontend.

One major challenge was ensuring consistency between the database schema,
API models, and frontend input handling. In particular, the handling of the
status attribute required careful debugging to ensure correct data flow
through all layers.

Another challenge was setting up the connection between FastAPI and PostgreSQL
in a Docker environment. Timing issues during container startup required
additional debugging.

Overall, the project provides a solid foundation for further extensions,
such as adding authentication, search functionality, or a more advanced UI.
% ============================================================
\section*{Sources and reuse}
% ============================================================

\addcontentsline{toc}{section}{Sources and reuse}

The project was implemented using concepts and code examples from the lecture
"Introduction to Database Management Systems" (THGA Bochum).

\begin{itemize}

  \item Database schema design and normalization concepts were based on lecture
        materials (DDL, DML, 3NF).

  \item The structure of the FastAPI application (routers, CRUD operations,
        database connection) was adapted from lecture examples and exercises.

  \item The Docker Compose setup was based on examples shown in the course
        (multi-container setup with PostgreSQL and API).

\end{itemize}

\end{document}
